\section{Synthesis of the results and consequences}
\label{nuclear:sintesis}
The generating structure preserves operations and history before its numerical realization. In realized transport, this conservation is formulated as an identity of forms. A partial reading distributes charge between the visible component and the complement; including both defines a recoverable transformation. Refinement prolongs the same identity without identifying different depths or replacing the register by its intensity.

\subsection{Conservation, reconstruction, and dimension}
The equality \(\mathcal U^*(b'\oplus b')\mathcal U=b\oplus b\) brings reading and memory together in a reversible operation. For compressions with a separate register, telescoping preserves the positive terminal component and the complete series of complements. For a fixed unitary holonomy, the terminal component is the projection onto its fixed-vector space; for a variable sequence, it is determined by the strong limit of the Gram operators.

Dodecaphasic memory allows the centred register to be recovered through two specified observations, with inverse norm \(9/4\). Incidence transports this recovery to the dimensional projector. Its exterior prolongation preserves each degree through Gram determinants and requires the inclusion of mixed components. Thus the relations between length, area, and volume are expressed through maps defined on the same register, with their respective invariants.

\subsection{Action and physical consequences}
Action, clock, and energy basis retain their types throughout composition. The electronic ratio \(m_ec^2/\hbar\) is expressed through the clock frequency and electronic character, preserving the basis-conversion factor. Maxwell--Dirac normalization removes the chosen dimensional scales and leaves the dimensionless coupling. Gravitational transport preserves the areal scale \(\hbar G/c^3\); the homogeneous bounce solution retains its conditions on matter, torsion, and cosmological constant.

The Noether relation is manifested when the transported observable is a charge of the declared action. Change of form adds a connection term to the Hamiltonian generator and preserves oriented action exactly. The holonomy of a return also preserves the information distinguishing phase return from return of the complete state.

\subsection{Memory and quantum spectrum}
Identity \eqref{nuclear:eq:espectro} determines the symplectic spectrum of the reduced covariance through the antilinear component of the memory operator. For the loop
\[
H_1=\begin{pmatrix}2&1\\1&1\end{pmatrix},
\qquad
W_1=\frac19\begin{pmatrix}13&3\\3&10\end{pmatrix},
\]
we obtain \(\det W_1=121/81\), symplectic eigenvalue \(11/9\), and purity \(9/11\). The spectrum of the reduced density operator is
\[
\lambda_j=\frac9{10}\,10^{-j},\qquad j=0,1,2,\ldots,
\]
and its entropy satisfies
\[
\frac{S}{k_B}=\frac{10}{9}\log10-\log9.
\]
These figures come from the transport and state fixed in the theorem. Spectral equivalence with a Gibbs density operator assigns a temperature to the chosen Hamiltonian; that equivalence retains its distinction from a dynamical thermalization process.

\subsection{Foundational content}
The articulation between number, incidence, and dynamics resides in the conservation of a common antecedent and in the operations relating its realizations. The article makes this articulation explicit through reconstructions, bounds, spectra, and actions. Generation of values and transport of symmetries retain the same provenance; their observable consequences are calculated subsequently with the reader and state that actually participate. This order provides a precise mathematical formulation of dynamical memory and its physical realizations.
